\documentclass[12pt]{article} % Corpo do texto em tamanho 12

% --- Configurações de Margem e Fonte ---
\usepackage[utf8]{inputenc}
\usepackage[T1]{fontenc}
\usepackage[a4paper, margin=2cm]{geometry} % Margem mínima de 2cm em todos os lados
\usepackage{times} % Fonte Times New Roman
\usepackage{xcolor}
\usepackage{listings}
\usepackage{graphicx}
\usepackage{fancyhdr}
\usepackage{hyperref}
\usepackage{titlesec} % Pacote para personalizar tamanhos de títulos

% --- Configuração de Tamanho de Títulos ---
% Títulos de seção (14pt ou 16pt). Definido aqui como 16pt para destaque.
\titleformat{\section}{\normalfont\Large\bfseries}{\thesection}{1em}{}
\titleformat{\subsection}{\normalfont\large\bfseries}{\thesubsection}{1em}{}

% --- Configuração de Blocos de Código (Fonte Monoespaçada) ---
\definecolor{codebg}{gray}{0.95} % Fundo cinza claro
\lstset{
    basicstyle=\ttfamily\small, % Fonte monoespaçada
    backgroundcolor=\color{codebg},
    breaklines=true,
    frame=single,
    rulecolor=\color{black!30},
    showstringspaces=false,
    captionpos=b % Legenda abaixo do bloco
}

% --- Cabeçalho e Rodapé ---
\pagestyle{fancy}
\fancyhf{}
\fancyhead[L]{Simple CTF - Lurian Lima} % Nome do desafio e autor no cabeçalho
\fancyhead[R]{Write-Up}
\fancyfoot[C]{\thepage} % Numeração de páginas obrigatória

\begin{document}

% --- Capa ---
\begin{titlepage}
    \centering
    \vspace*{2cm}
    {\Huge \textbf{Write-Up de Segurança Ofensiva} \par}
    \vspace{1.5cm}
    {\Large \textbf{Desafio:} Simple CTF \par}
    {\large \textbf{Plataforma:} Try Hack Me \par}
    {\large \textbf{Autor:} Lurian Lima \par}
    \vspace{2cm}
    {\large \textbf{Data de Resolução:} \today \par}
    \vfill
\end{titlepage}

% --- Sumário ---
\tableofcontents
\newpage

% ---------------
% Resultados da enumeração. Use o comando abaixo para inserir screenshots:
% \begin{figure}[H]
%     \centering
%     \includegraphics[width=0.8\textwidth]{screenshot_nmap.png}
%     \caption{Legenda descritiva da imagem inserida imediatamente após o texto.}
% \end{figure}

% \begin{lstlisting}[caption={Comando de exploração utilizado}]
% # Exemplo de bloco de código com fundo destacado
% nmap -sV -sC -oN scan.txt 10.10.10.10
% \end{lstlisting}
% ---------------

% --- 2.3 Visão Geral ---
\section{Visão Geral}
% Parágrafo descrevendo o vetor principal de ataque sem spoiler detalhado. Deve responder: o que é o desafio,
% qual a superfície de ataque explorada e qual foi o caminho geral até o objetivo.

% --- 2.4 Reconhecimento ---
\section{Reconhecimento}
% Resultados da enumeração externa: portas abertas, serviços identificados, versões relevantes, tecnologias
% detectadas. Incluir comandos usados e os trechos de output mais relevantes. Não incluir informações obtidas
% após o foothold inicial.

% --- 2.5 Exploração ---
\section{Exploração}
% Descrição técnica do vetor de entrada. Deve incluir o raciocínio que levou à identificação da vulnerabilidade,
% os payloads ou scripts utilizados (comentados), exploits públicos utilizados e os outputs que confirmam o
% acesso. Cada passo deve ser justificado, não apenas listado.

% --- 2.6 Movimentação Lateral (Se aplicável) ---
\section{Movimentação Lateral}
% Aplicável quando há múltiplos usuários ou pivoting entre serviços. Descrever a escalada horizontal com o
% mesmo padrão da seção anterior.

% --- 2.7 Escalada de Privilégios ---
\section{Escalada de Privilégios}
% Descrever o vetor de privesc identificado, o processo de enumeração local que o revelou (linpeas, sudo -l,
% capabilities, etc.) e a exploração realizada. Incluir o comando final que comprova o acesso root/SYSTEM

% --- 2.8 Conclusão ---
\section{Conclusão}
% Resumo técnico de 3 a 5 linhas sintetizando a cadeia de exploração completa, do reconhecimento à flag.

% --- 2.9 Referências ---
\section{Referências}
% Espaço destinado às fontes utilizadas para a exploração: links dos repositórios dos exploits públicos e links
% das CVEs correspondentes.

\end{document}